\documentclass[../main.tex]{subfiles}

\begin{document}

\chapter{完备度量空间与函数空间}\label{chapter:ch07}

度量空间的完备性概念, 读者可能早就知道了, 它是分析学各方面的基础. 虽然完备性只是一种度量性质, 而非拓扑性质, 但仍然有许多涉及完备度量空间的定理具有拓扑的色彩. 本章中, 我们将研究一些重要的完备度量空间, 并且证明若干相关的定理.

完备度量空间的最常见的例子是通常度量下的欧氏空间. 另一个重要例子是把一个空间 $X$ 映入一个度量空间 $Y$ 的所有连续函数构成的集合 $\mathcal{C}(X,Y)$. 这个集合上有一个度量称之为一致度量, 它和第 20 节中给 $\mathbb{R}^J$ 定义的那种一致度量类似. 如果 $Y$ 是一个完备度量空间, 那么 $\mathcal{C}(X,Y)$ 在一致度量下也是完备的. 这一点我们将在第 43 节中予以证明. 作为一个应用, 我们在第 44 节中构造了著名的充满空间的 Peano 曲线.

涉及完备度量空间的拓扑特性的定理之一便是关于空间的紧致性与完备性关系的定理, 我们在第 45 节中证明它. 这个定理的一个直接的推论是函数空间 $\mathcal{C}(X,\mathbb{R}^n)$ 中关于紧致子空间的一个定理, 这就是著名的 Ascoli 定理的一种经典形式.

在函数空间 $\mathcal{C}(X,Y)$ 上, 除了由一致度量所诱导的拓扑之外, 还有其他有用的拓扑. 在第 46 节中, 我们将对其中的某些拓扑予以研究, 并在第 47 节中证明 Ascoli 定理的一般形式.

\subfile{Chapter-07/section-43}
\subfile{Chapter-07/section-44}
\subfile{Chapter-07/section-45}
\subfile{Chapter-07/section-46}
\subfile{Chapter-07/section-47}

\end{document}